\documentclass[11pt,a4paper]{article}
\usepackage[utf8]{inputenc}
\usepackage[T1]{fontenc}
\usepackage[spanish,es-noquoting]{babel}
\usepackage[margin=2.5cm]{geometry}
\usepackage{longtable}
\usepackage{booktabs}
\usepackage{enumitem}
\usepackage{fancyhdr}
\usepackage[hidelinks]{hyperref}

\pagestyle{fancy}
\fancyhf{}
\lhead{\small Protocolo del componente empírico --- SIGA}
\rhead{\small\thepage}
\renewcommand{\headrulewidth}{0.4pt}

\title{\textbf{Protocolo del componente empírico}\\[0.3em]
\large Proyecto SIGA --- Sistema Inteligente de Gestión de Aulas}
\author{Ingeniería de Requerimientos [20303] --- Equipo FGMMN\\
Facultad de Ciencias de la Computación, UTEQ}
\date{Versión 2.0 --- 3 de agosto de 2026}

\begin{document}
\maketitle

\section*{Enfoque}
Comparación de calidad de requisitos funcionales elicitados por un equipo humano
frente a los generados por un modelo de lenguaje grande (LLM) a partir del mismo
material fuente.

\section{Marco PICOC}
\begin{longtable}{p{3.2cm}p{11cm}}
\toprule
\textbf{Elemento} & \textbf{Definición} \\
\midrule
\endhead
Población & Requisitos funcionales del sistema SIGA (26 generados por LLM, 25 elicitados por el equipo humano). \\
Intervención & Elicitación asistida por modelo de lenguaje grande, a partir del corpus completo de 11 entrevistas anonimizadas. \\
Comparación & Elicitación tradicional por el equipo humano (entrevistas y análisis manual). \\
Resultado & Puntuación de calidad en 5 dimensiones: completitud, ausencia de ambigüedad, verificabilidad, corrección respecto de la fuente y consistencia interna; evaluada por al menos 3 jueces ciegos. \\
Contexto & Proyecto académico de Ingeniería de Requerimientos, dominio de gestión inteligente de aulas, FCCDD, UTEQ. \\
\bottomrule
\end{longtable}

\section{Pregunta de investigación}
¿En qué dimensiones de calidad los requisitos funcionales elicitados por un equipo
humano difieren de los generados por un modelo de lenguaje grande a partir del mismo
material fuente?

\section{Hipótesis}
\begin{itemize}[leftmargin=*]
  \item \textbf{H\textsubscript{0}:} No hay diferencia significativa en la puntuación
        media de calidad entre el Conjunto A (LLM) y el Conjunto B (humano) en
        ninguna de las 5 dimensiones evaluadas.
  \item \textbf{H\textsubscript{1}:} Existe una diferencia significativa en la
        puntuación media de calidad entre el Conjunto A y el Conjunto B en al menos
        una dimensión.
\end{itemize}

\section{Variables}
\begin{longtable}{p{3.2cm}p{11cm}}
\toprule
\textbf{Tipo} & \textbf{Variable} \\
\midrule
\endhead
Independiente & Origen del requisito (Humano / LLM). \\
Dependientes & Puntuación de 1 a 5 en cada una de las 5 dimensiones de la rúbrica. \\
Control & Mismo material fuente para ambos conjuntos; los mismos jueces evalúan ambos conjuntos; orden aleatorizado y ciego. \\
\bottomrule
\end{longtable}

\section{Participantes (jueces evaluadores)}
\textbf{Requisito de independencia:} ninguno de los jueces puede ser una de las
personas entrevistadas en las 11 transcripciones que sirvieron como material fuente
(DOC-01 a DOC-04, COORD-01 a COORD-03, CONS-01 a CONS-04), para preservar la ceguera
del diseño. Se exige un mínimo de 3 jueces; se recomiendan 5 cuando sea posible,
conforme a la Sección 5.1 de la guía.

Cada juez firmó un consentimiento informado específico, archivado en
\texttt{06\_Experimento/instrumentos/}.

\section{Diseño}
Cuasi-experimento apareado. Cada juez evalúa los 51 ítems del paquete ciego
(\texttt{Paquete\_Evaluacion\_Ciega\_Jueces.md}), sin conocer el origen de cada uno,
usando la \texttt{Hoja\_Puntuacion\_JUEZ.csv}. El orden fue aleatorizado con semilla
fija (20260802) para garantizar reproducibilidad.

\section{Plan de análisis estadístico}
\begin{enumerate}[leftmargin=*]
  \item Consolidar las hojas de puntuación de los jueces.
  \item Calcular el acuerdo inter-evaluador: $\kappa$ de Cohen por par de jueces y
        $\kappa$ de Fleiss para el conjunto completo.
  \item Prueba de normalidad de Shapiro-Wilk sobre las puntuaciones de cada conjunto,
        por dimensión.
  \item Si los datos son normales, prueba $t$ para muestras apareadas comparando la
        media de cada juez en el Conjunto A frente al Conjunto B. Si no lo son,
        prueba de Wilcoxon.
  \item Tamaño del efecto: $d$ de Cohen en el caso paramétrico, $\delta$ de Cliff en
        el no paramétrico.
  \item Cálculo de potencia estadística previo, con $\alpha = 0{,}05$ y potencia
        deseada $1-\beta = 0{,}80$.
\end{enumerate}

Los scripts que reproducen estas tablas a partir de los datos crudos se encuentran en
\texttt{06\_Experimento/scripts\_analisis/}.

\section{Amenazas a la validez}
\begin{itemize}[leftmargin=*]
  \item \textbf{Validez de constructo:} el modelo usado para generar el Conjunto A
        tuvo exposición previa parcial al Conjunto B dentro de la misma cuenta de
        chat. La limitación se documenta en detalle en
        \texttt{prompt\_llm\_conjunto\_A.md}.
  \item \textbf{Validez interna:} los jueces deben ser verificablemente independientes
        de las 11 entrevistas fuente; su selección queda documentada.
  \item \textbf{Validez externa:} los resultados son específicos del dominio SIGA y de
        un modelo de LLM particular; no se generalizan automáticamente a otros
        dominios ni a otros modelos.
  \item \textbf{Validez de conclusión:} el tamaño de muestra (51 ítems, 3 jueces) es
        modesto. Se recomienda ampliar el número de jueces hacia la Entrega 4 (2B) si
        el tiempo lo permite.
\end{itemize}

\section{Registro previo del protocolo}
El protocolo está registrado en el Marco de Ciencia Abierta (OSF):

\begin{center}
\url{https://osf.io/3nruf/}
\end{center}

\textbf{Estado actual:} \emph{Pending registration approval}. El registro fue
enviado antes de la distribución del paquete ciego a los jueces, conforme al
gatekeeper G6 de la guía. La aprobación por parte de OSF es un trámite administrativo
posterior al envío y no altera la fecha de registro, que es la que otorga la
precedencia temporal.

\section{Disponibilidad de datos y materiales}
El conjunto de datos anonimizado está depositado en Zenodo bajo licencia Creative
Commons Atribución 4.0 Internacional (CC BY 4.0):

\begin{center}
\url{https://doi.org/10.5281/zenodo.21774351}
\end{center}

Los datos identificables (grabaciones de audio y video, consentimientos con cédula y
firma visibles) permanecen en la zona restringida del repositorio, cifrada con
AES-256, conforme a la Sección 4.1 de la guía y a la Ley Orgánica de Protección de
Datos Personales del Ecuador.

\end{document}
